\subsection{Experimental Setup}
\label{sec:experiments}

The executable study defines the four frozen tasks in Band A of
Figure~\ref{fig:pipeline} rather than comparing errors across unlike
resolutions. Their complete contexts, outputs, evaluation horizons, and origin
counts are retained in the Supplementary Material. Within a task, all models
share sites, origins, targets, fitting cuts, and post-processing. Across tasks,
MAE and RMSE retain W/m$^2$ units but refer to different temporal aggregates;
they are not pooled.

\subsubsection{Common matched protocol}

The daily and monthly experiments contain three deterministic references
(persistence, annual seasonal na\"ive, and calendar climatology) and four
learned models (XGBoost, LSTM, TimesNet, and DilatedRNN). Every learned
model receives only causal GHI and site identity. Neural models use the five
seeds 11, 23, 42, 67, and 89, and XGBoost is likewise refitted under each of
those seeds because row and column subsampling are stochastic; both seed-level
forecasts and the arithmetic mean of every learned model are retained.
Deterministic references are evaluated once. The primary model comparison uses
macro-MAE after computing
MAE separately at each site. RMSE, nRMSE, and $R^2$, as well as cumulative and
exact-lead scopes, remain mandatory complementary outputs.

For selection, neural networks are initialized under each seed, fitted to the
adjustment windows, and monitored on the 2023 windows using normalized MSE.
Early stopping chooses an epoch count for each model--seed pair. The network is
then discarded, reinitialized under the same seed, and fitted to the pre-2024
refit windows for exactly that number of epochs. Thus, epoch selection
minimizes validation MSE in site-normalized units, while the reported primary
ranking uses test MAE in physical units. XGBoost uses its frozen tree
budget in separate adjustment and refit instances. Test forecasts are generated
only by the refitted instances.

Band B of Figure~\ref{fig:pipeline} summarizes this selection--refit chronology;
Band A shows the four direct-output contracts that remain fixed during testing.

Within each partition, forecasts follow the rolling-origin protocol of
Section~\ref{sec:methodology}: the official hourly task advances one local-
midnight origin per day, the daily task advances one origin per day, and the
monthly tasks advance one origin per civil month. Parameters remain fixed in
the test partition while each context incorporates only observations that have
become available before its first target.

\subsubsection{Hourly protocol and DilatedRNN extension}

The official hourly task uses 336 past hours (14 days) and produces the next
72 hours directly. The reported checkpoints are 24, 48, and 72 hours; they are
prefixes and exact leads extracted from the same 72-hour output, not separately
trained models. Adjustment origins run from 15 January 2019 through
29 December 2022 (1,445 per site), and validation origins run from
1 January through 29 December 2023 (363 per site). The pre-test refit combines
origins from 15 January 2019 through 29 December 2023 (1,810 per site). The
retrospective test contains 363 origins from 1 January through
28 December 2024. The last 72-hour window ends at 23:00 on 30 December in the
fixed local-time index; corresponding UTC timestamps can extend into
31 December for western locations.

The official artifact fixes seed 42 and contains persistence, annual seasonal
na\"ive, XGBoost, LSTM, and TimesNet predictions. The extension trains only a
direct DilatedRNN with seed 42 under the same 336-to-72 contract. Its forecasts
are merged one-to-one with the read-only official artifact by site, forecast
origin, target timestamp, and lead. Target equality is checked during the
one-to-one merge, and the hashes of all official inputs are rechecked before
the extension is marked complete. Thus the hourly extension does not silently
retrain or overwrite an official comparator.

\subsubsection{Daily protocol}

The daily primary task uses 365 daily means to predict a 30-day vector. Metrics
are retained at 1, 3, 7, 14, and 30 days. Adjustment origins extend from
1 January 2020 through 2 December 2022, with targets ending no later than
31 December 2022 (1,067 origins per site). Validation uses
1 January--2 December 2023 (336 origins), and refit uses
1 January 2020--2 December 2023 (1,432 origins). The 2024 test uses
1 January--2 December and contains 337 complete origins per site because 2024
is a leap year.

\subsubsection{Monthly primary and exploratory protocols}

Civil-month means are calculated from the audited daily series. The primary
monthly task uses 12 months to predict the next month. It has 36 adjustment
origins in 2020--2022, 12 validation origins in 2023, 48 refit origins in
2020--2023, and 12 test origins in 2024, for every site.

The exploratory monthly task uses the same 12-month context and directly
predicts six months. Its declared checkpoints are months 1, 3, and 6.
The final origin in each partition must fall no later than July so that all
six targets remain inside its final year. Adjustment therefore contains
31 consecutive monthly origins from January 2020 through July 2022, validation
contains seven from January through July 2023, refit contains 43 from January
2020 through July 2023, and test contains seven from January through July 2024.
Earlier complete years in adjustment and refit retain all twelve origins.
This task is explicitly labelled exploratory because each site contributes
only seven test origins, each with a six-month output.

For the six-month task, exact leads are evaluated on different target-month
sets: lead 1 covers January--July, whereas lead 6 covers June--December. Its
horizon profile therefore combines lead-time and seasonal effects and is not
interpreted as an isolated degradation curve.

\subsubsection{Frozen optimization settings}

Table~\ref{tab:hyperparameters} gives the compact configuration needed to
interpret the main comparison. The Supplementary Material records the
full temporal boundaries, model settings, baseline definitions, artifact
contract, and external-context integrity hashes.

\begin{table}[!t]
\centering
\caption{Compact fixed configuration of the matched implementation.}
\label{tab:hyperparameters}
\footnotesize
\setlength{\tabcolsep}{3pt}
\renewcommand{\arraystretch}{1.12}
\begin{tabular}{@{}
  >{\raggedright\arraybackslash}p{0.27\linewidth}
  >{\raggedright\arraybackslash}p{0.66\linewidth}@{}}
\toprule
\textbf{Component} & \textbf{Fixed setting} \\
\midrule
Matched input & Causal GHI sequence plus fixed site identifier only \\
Regression/output & Continuous GHI; direct vector of task-specific length \\
Neural optimization & Adam; MSE loss; learning rate $10^{-3}$;
weight decay $10^{-5}$ \\
Batch size & 128 for hourly/daily; 32 for monthly \\
Epoch selection & Hourly/daily: maximum 30, patience 5; monthly:
maximum 300, patience 30 \\
Seeds & Daily/monthly: 11, 23, 42, 67, 89; official hourly and
DilatedRNN extension: 42 \\
TimesNet & $d_{\mathrm{model}}=8$, $d_{\mathrm{ff}}=16$, one TimesBlock,
top-$k=3$, two kernels, dropout 0.1 \\
DilatedRNN & Dilations $(1,2,4)$; 16 units/layer; dense width 16;
dropout 0 \\
LSTM & One layer; hidden width 32; direct output head \\
Site embedding & Width 8 for all neural models \\
XGBoost & 150 trees; depth 3; learning rate 0.05; row and column
subsampling 0.8; histogram tree method \\
Post-processing & Zero floor for all tasks; hourly values also set to zero
when solar elevation is nonpositive \\
\bottomrule
\end{tabular}
\end{table}


\subsubsection{Model capacity and computational evidence}

To make differences in model capacity explicit,
Table~\ref{tab:model_capacity} gives the exact number of trainable scalar
parameters in each neural model. The counts are obtained from the frozen module
dimensions, including biases and the ten-site embedding, rather than inferred
from checkpoint file sizes. TimesNet has a task-dependent direct temporal
projection, so its parameter count grows with both context and output length.
Dilated recurrence changes the temporal connections without adding separate
weights for each skipped position. XGBoost is instead constrained by 150 trees
of maximum depth three and is not assigned an artificial neural-parameter
equivalent.

\input{tables/model_capacity}

Training and inference times were not measured consistently across the
completed runs. The computational comparison therefore reports parameter
counts and the common one-thread policy; runtime comparisons remain for
future evaluation.
